\chapter{Creatine Kinase (CK)}

\section*{What Is This Test?}
Creatine kinase (CK), also called creatine phosphokinase (CPK), is an enzyme that catalyzes the reversible transfer of a high-energy phosphate group from creatine phosphate to ADP, regenerating ATP from ADP during periods of high energy demand: Creatine phosphate + ADP $\leftrightarrow$ Creatine + ATP. This reaction maintains ATP levels in tissues with high and fluctuating energy demands, primarily skeletal muscle, cardiac muscle, and brain. CK exists as three isoenzymes, each a dimer composed of two subunits: M (muscle) and B (brain). CK-MM (CK-3) is found predominantly in skeletal muscle (>98\% of muscle CK), CK-MB (CK-2) is found primarily in cardiac muscle (approximately 20--30\% of cardiac CK, with the remainder being CK-MM), and CK-BB (CK-1) is found mainly in brain and smooth muscle (gastrointestinal tract, uterus, bladder). In healthy adults, serum total CK is derived almost exclusively from skeletal muscle (CK-MM), with very low CK-MB (<5\% of total, typically <5 ng/mL or <2--3\% as a percentage). CK rises in the blood 4--8 hours after muscle injury, peaks at 12--24 hours, and returns to baseline in 48--72 hours (half-life approximately 15 hours). Total CK is the most sensitive enzyme marker for skeletal muscle injury. For cardiac injury, troponin I or T has replaced CK-MB as the gold standard biomarker for myocardial infarction; however, CK-MB is still used in some contexts, such as estimating infarct size, detecting reinfarction (because CK-MB normalizes more quickly than troponin), and confirming periprocedural MI after cardiac surgery or PCI.

\section*{Alternative Names}
CPK; Creatine Phosphokinase; Total CK; CK-MB; CK-MM; CK-BB; Serum Creatine Kinase; CK Isoenzymes; CK-MB Mass

\section*{Principle}
Total CK is measured by an enzymatic kinetic method based on the reverse reaction (creatine phosphate formation is favored at pH 6.8): Creatine phosphate + ADP $\rightarrow$ Creatine + ATP (CK). The ATP produced is coupled to a hexokinase/glucose-6-phosphate dehydrogenase (G6PD) indicator system: ATP + Glucose $\rightarrow$ ADP + Glucose-6-phosphate (hexokinase). Glucose-6-phosphate + NADP\textsuperscript{+} $\rightarrow$ 6-Phosphogluconate + NADPH + H\textsuperscript{+} (G6PD). The rate of increase in absorbance at 340 nm (as NADPH is produced) is proportional to CK activity. The IFCC standardized method uses N-acetylcysteine to reactivate oxidized CK sulfhydryl groups. Results are reported in IU/L at 37 degrees C. CK-MB is measured by immunoassay (mass assay) using monoclonal antibodies specific to the CK-MB dimer (reported in ng/mL) or by activity inhibition using an anti-CK-M antibody that selectively inhibits the M subunit (reported in IU/L). The mass assay is more sensitive and specific than the activity assay and is the standard method. The CK-MB relative index (RI) = CK-MB mass / Total CK $\times$ 100\%. RI >2.5--3\% suggests cardiac origin. Results must be interpreted with extreme caution in patients with concurrent skeletal muscle injury. This test is NOT performed by hematology analyzers.

\section*{History}
Creatine was discovered by Michel Chevreul in 1832. Creatine kinase was discovered by Lohmann in 1934. The CK isoenzymes were separated by electrophoresis in the 1960s. CK-MB was established as a cardiac biomarker in the 1970s and became the gold standard for MI diagnosis for three decades. The landmark WHO criteria for MI diagnosis (1979) used CK-MB as a core component. Troponin assays were introduced in the 1990s and largely replaced CK-MB due to superior sensitivity and specificity. CK-MB remains useful for estimating infarct timing (rises 4--6 hours, peaks 12--24 hours, normalizes 48--72 hours) and detecting reinfarction (troponin stays elevated for 7--14 days, making reinfarction diagnosis difficult).

\section*{Why This Test Is Ordered}
\begin{itemize}
  \item Diagnosis and monitoring of rhabdomyolysis: CK >5$\times$ ULN (>1,000--2,000 IU/L or >5--10$\times$ ULN) with appropriate clinical context (trauma, crush injury, prolonged immobilization, seizures, drugs, exertion) defines rhabdomyolysis. CK typically exceeds 5,000--10,000 IU/L in clinically significant rhabdomyolysis and may exceed 100,000 IU/L in severe cases
  \item Detection of statin-induced myopathy: CK >10$\times$ ULN with myalgia, weakness, or myoglobinuria defines statin-associated rhabdomyolysis. Mild CK elevation (3--5$\times$ ULN) with muscle symptoms suggests myositis
  \item Evaluation of suspected inflammatory myopathies: polymyositis, dermatomyositis, necrotizing autoimmune myopathy. CK often elevated 5--50$\times$ ULN
  \item Adjunctive diagnosis of myocardial infarction: CK-MB mass with relative index; used primarily for reinfarction diagnosis or when troponin is equivocal/unavailable
  \item Detection of muscular dystrophy: CK markedly elevated in Duchenne muscular dystrophy (often >10,000--30,000 IU/L from birth, declines with age as muscle mass is lost) and Becker muscular dystrophy
  \item Evaluation of unexplained myalgia, weakness, dark urine (myoglobinuria), or elevated AST with normal ALT
  \item Screening for malignant hyperthermia susceptibility (elevated resting CK in approximately 50\% of susceptible individuals)
  \item Monitoring of crush injury and compartment syndrome in trauma and disaster settings
\end{itemize}

\section*{Is This a Primary or Secondary Test}
Total CK is a secondary test for specific clinical scenarios (suspected rhabdomyolysis, myopathy, muscular dystrophy). For myocardial infarction, high-sensitivity troponin I or T is the current gold standard, and CK-MB is a secondary adjunct when needed.

\section*{How This Test Is Performed}
A healthcare professional draws blood from a vein into a serum separator tube (gold-top) or lithium heparin tube (green-top). Important: EDTA (purple-top) tubes must NOT be used because EDTA chelates magnesium, which is a required cofactor for CK activity, producing falsely low results. Prolonged tourniquet time should be avoided. The sample is centrifuged and analyzed on an automated chemistry analyzer. Results are available within 1--2 hours. CK is light-sensitive to some degree---samples should be protected from prolonged light exposure. Moderate hemolysis can falsely elevate CK because red blood cells contain adenylate kinase, which interferes with the CK assay by producing ATP from ADP.

\section*{What Preparation Is Needed}
No special preparation is required. However, several factors significantly affect CK: (1) Vigorous exercise (weightlifting, marathon running, CrossFit) can elevate CK by 5--25$\times$ ULN, peaking at 24--48 hours post-exercise and normalizing over 7--10 days. This is a normal physiological response to muscle microtrauma (exertional rhabdomyolysis) and should not be confused with pathologic rhabdomyolysis unless there is myoglobinuria, renal failure, or CK >50,000 IU/L. (2) Intramuscular (IM) injections cause mild CK elevation (2--5$\times$ ULN) from local muscle trauma. (3) Recent surgery, trauma, or prolonged immobilization (e.g., floor after a fall with inability to get up for hours) causes significant CK elevation. (4) Medications: statins (especially high-dose atorvastatin and simvastatin), fibrates (gemfibrozil), colchicine, hydroxychloroquine, zidovudine (mitochondrial myopathy), succinylcholine, antipsychotics (neuroleptic malignant syndrome), propofol (propofol infusion syndrome). (5) Alcohol and cocaine use elevate CK. (6) Intramuscular injections of certain medications (ketorolac, ceftriaxone, haloperidol) can cause CK elevation. (7) Seizures, particularly generalized tonic-clonic, elevate CK dramatically (>3,000 IU/L within 24 hours). (8) Hypothyroidism causes elevated CK with CK-MM predominance (from reduced clearance and myopathy). (9) Race: black individuals have CK levels approximately 50--100\% higher on average than white individuals due to greater muscle mass and possible genetic differences.

\section*{Contraindications and Precautions}
No absolute contraindications. Critical considerations: (1) CK rises in virtually any form of muscle injury or strenuous exercise; clinical history is essential to differentiate pathological from physiological elevation. (2) CK levels are higher in males, in individuals with greater muscle mass, and in black individuals. Reference ranges are sex and race-specific. (3) CK levels in hypothyroidism can be substantially elevated (5--50$\times$ ULN) without muscle symptoms. TSH should be checked. (4) Macro-CK (macroenzyme type 1: CK-BB complexed with IgG; type 2: mitochondrial CK oligomers) can cause persistently elevated total CK. Suggested by atypical isoenzyme pattern on electrophoresis, and by PEG precipitation. (5) Neuroleptic malignant syndrome (NMS) presents with CK >1,000 IU/L (often 10,000--50,000), hyperthermia, rigidity, altered mental status, and autonomic instability after antipsychotic administration. (6) Malignant hyperthermia: CK >20,000 IU/L with hyperthermia, tachycardia, hypercapnia, and muscle rigidity after exposure to succinylcholine or volatile anesthetics. (7) CK is normal in disorders of neuromuscular transmission (myasthenia gravis, Lambert-Eaton syndrome) and in most neurogenic disorders (ALS, neuropathy) unless there is secondary muscle involvement.

\section*{Risks}
Minimal risk from venipuncture. Mild pain or bruising at site resolves in 1--2 days.

\section*{What Happens After the Test}
Results available within 1--2 hours. Total CK interpretation: (1) CK 200--1,000 IU/L: mild elevation; consider recent exercise, IM injection, trauma, early myopathy, hypothyroidism, statin effect. (2) CK 1,000--5,000 IU/L: moderate elevation; consider rhabdomyolysis, myositis, prolonged immobilization, seizure, NMS, malignant hyperthermia, severe statin myopathy. (3) CK 5,000--50,000 IU/L: rhabdomyolysis; risk of acute kidney injury from myoglobin tubular toxicity. CK >5,000 IU/L with AKI is diagnostic of rhabdomyolysis. (4) CK >100,000 IU/L: massive rhabdomyolysis; usually from crush injury, severe NMS/malignant hyperthermia, status epilepticus, or severe exertional rhabdomyolysis. Aggressive IV fluid resuscitation (200--300 mL/hour of normal saline, targeting urine output 200--300 mL/hour) is indicated to prevent myoglobinuric AKI. Rhabdomyolysis complications: hyperkalemia (life-threatening from muscle cell potassium release), hypocalcemia (calcium deposition in damaged muscle, often followed by rebound hypercalcemia during recovery), hyperphosphatemia, hyperuricemia, compartment syndrome, and AKI (myoglobin is directly nephrotoxic and precipitates with Tamm-Horsfall protein in acidic urine).

\section*{Understanding Results}

\subsection*{Below Normal}
\begin{itemize}
  \item Low CK is not clinically significant. It may reflect low muscle mass (elderly, cachexia, prolonged immobilization, steroid myopathy, disuse atrophy) or early pregnancy.
  \item In Duchenne muscular dystrophy, CK is markedly elevated in early childhood (often >10,000 IU/L) but declines with age as muscle mass is lost; a low CK in a 15-year-old with DMD reflects end-stage disease, not improvement.
  \item In myasthenia gravis, CK is normal.
  \item Low CK does not exclude cardiac injury (troponin is the preferred cardiac biomarker).
\end{itemize}

\subsection*{Above Normal}
\begin{itemize}
  \item \textbf{Rhabdomyolysis:} CK >5--10$\times$ ULN (typically >1,000--5,000 IU/L; often >10,000 IU/L). Causes: trauma (crush injury, motor vehicle accident), prolonged immobilization (lying on floor after stroke, fall, or drug overdose for hours), extreme exertion (marathon, CrossFit, military training, status epilepticus), drugs/toxins (statins, fibrates, alcohol, cocaine, heroin, amphetamines, antipsychotics causing NMS, succinylcholine causing malignant hyperthermia, colchicine, propofol), infections (viral myositis---influenza, Coxsackie, HIV), electrolyte disorders (severe hypokalemia, hypophosphatemia), ischemia (compartment syndrome, arterial thrombosis, prolonged tourniquet), metabolic myopathies (McArdle disease---glycogen phosphorylase deficiency; carnitine palmitoyltransferase II deficiency), and heat stroke.
  \item \textbf{Myocardial infarction:} Total CK may be mildly elevated if skeletal muscle injury coexists. CK-MB mass >5 ng/mL with CK-MB relative index >2.5--3\% suggests cardiac source. CK-MB rises 4--6 hours after MI, peaks at 12--24 hours, normalizes at 48--72 hours. This shorter normalization window (vs. troponin's 7--14 days) makes CK-MB useful for detecting reinfarction.
  \item \textbf{Inflammatory myopathies:} Polymyositis (CK 5--50$\times$ ULN, symmetric proximal weakness, anti-Jo-1 often positive), dermatomyositis (similar, with Gottron papules, heliotrope rash, shawl sign, V-sign), necrotizing autoimmune myopathy (often associated with statin exposure or anti-SRP or anti-HMGCR antibodies; CK often >10--50$\times$ ULN), inclusion body myositis (CK normal or mildly elevated <10$\times$ ULN, distal and proximal weakness, older men).
  \item \textbf{Muscular dystrophies:} Duchenne MD (CK >10,000--30,000 IU/L from birth, X-linked, progressive weakness, calf pseudohypertrophy, loss of ambulation by age 12). Becker MD (milder, CK 5,000--20,000 IU/L). Limb-girdle MD (CK variable, often 500--10,000 IU/L). Myotonic dystrophy types 1 and 2 (CK normal or mildly elevated).
  \item \textbf{Non-pathological elevation:} Vigorous exercise (CK peaks at 24--48 hours, up to 5--25$\times$ ULN, resolves within 7--10 days). Recent intramuscular injections. Surgery. Increased muscle mass (bodybuilders). Black race (CK ~50--100\% higher). Hypothyroidism (2--50$\times$ ULN). Medications (statins, fibrates, colchicine, antipsychotics).
\end{itemize}

\section*{Normal Results}
\begin{itemize}
  \item \textbf{Total CK:} Males: 55--170 IU/L (may be up to 200--300 IU/L in muscular/black individuals); Females: 30--145 IU/L
  \item \textbf{Total CK (black males):} Up to 400--500 IU/L may be normal
  \item \textbf{CK-MB mass:} <5 ng/mL (or <3.3--5.0 ng/mL depending on assay)
  \item \textbf{CK-MB relative index (RI):} <2.5--3\% of total CK
  \item \textbf{CK-MB activity:} 0--5 IU/L (<6\% of total CK activity)
  \item \textbf{CK peak post-exercise:} 24--48 hours post-exercise; normalize 7--10 days
  \item \textbf{CK half-life:} 15 hours
\end{itemize}

\section*{Follow-up Tests}
\begin{itemize}
  \item \textbf{CK-MB mass and relative index:} If myocardial origin is suspected.
  \item \textbf{High-sensitivity cardiac troponin I or T:} Gold standard for myocardial infarction.
  \item \textbf{Renal function (BUN, creatinine) and urinalysis:} Monitor for AKI and myoglobinuria in rhabdomyolysis. Urine dipstick positive for blood (heme) without RBCs on microscopy suggests myoglobinuria.
  \item \textbf{Serum electrolytes (potassium, calcium, phosphate):} Hyperkalemia (life-threatening), hypocalcemia, hyperphosphatemia in rhabdomyolysis.
  \item \textbf{Complete metabolic panel:} AST elevated from muscle (CK > AST); uric acid elevated.
  \item \textbf{TSH:} To evaluate for hypothyroidism.
  \item \textbf{CK isoenzyme electrophoresis:} Resolves unexplained CK elevation; identifies macro-CK.
  \item \textbf{EMG and muscle biopsy:} For suspected myopathy or muscular dystrophy.
  \item \textbf{Genetic testing:} For Duchenne/Becker (dystrophin gene), limb-girdle MD, metabolic myopathies.
\end{itemize}

\section*{References}
\begin{enumerate}
  \item Burtis CA, Ashwood ER, Bruns DE. Tietz Textbook of Clinical Chemistry and Molecular Diagnostics. 7th ed. Elsevier; 2023.
  \item Wallach J. Interpretation of Diagnostic Tests. 10th ed. Wolters Kluwer; 2022.
  \item Bosch X, Poch E, Grau JM. Rhabdomyolysis and acute kidney injury. N Engl J Med. 2009;361(1):62--72.
  \item Zimmerman JL, Shen MC. Rhabdomyolysis. Chest. 2013;144(3):1058--1065.
  \item Dalakas MC. Inflammatory muscle diseases. N Engl J Med. 2015;372(18):1734--1747.
  \item Thompson PD, Clarkson P, Karas RH. Statin-associated myopathy. JAMA. 2003;289(13):1681--1690.
  \item Morrow DA, Cannon CP, Jesse RL, et al. National Academy of Clinical Biochemistry Laboratory Medicine Practice Guidelines: clinical characteristics and utilization of biochemical markers in acute coronary syndromes. Circulation. 2007;115(13):e356--e375.
  \item Panteghini M. Recommendations for the use of cardiac troponin and CK-MB mass in acute coronary syndromes. Clin Chem Lab Med. 2006;44(5):550--556.
\end{enumerate}

\clearpage
